\newpage
\section{Acta Número 14}
\label{sec:acta14}

\noindent \textbf{Fecha de la sesión:}\par
Lunes, 20 de Abril de 2026

\vspace*{0.5cm}
\noindent \textbf{Datos de la Sesión:}
\vspace{-0.2cm}
\begin{itemize}[noitemsep]
    \item \textbf{Modalidad:} Virtual - Google Meet.
    \item \textbf{Hora de Inicio:} 18:00
    \item \textbf{Hora de Fin:} 20:00
\end{itemize}

\vspace*{0.5cm}
\noindent \textbf{Orden del Día:}
\vspace{-0.2cm}
\begin{itemize}[noitemsep]
    \item Asistencia.
    \item Lanzamiento (\textit{Kick-off}) de la fase de Aseguramiento de Calidad (QA) del Incremento.
    \item Asignación de roles de QA por células de desarrollo.
    \item Actualización del \textit{Definition of Done} (DoD): Integración de Casos de Prueba y Reportes de Errores.
    \item Revisión técnica de impedimentos (Backend, Frontend y Base de Datos).
    \item Gestión del \textit{Sprint Backlog}: Análisis de bloqueos y \textit{Carry Over}.
    \item Alineación técnica sobre la gestión del repositorio de documentación.
    \item Firmas y Aprobación de los Presentes.
\end{itemize}

\vspace*{0.5cm}
\noindent \textbf{Socios Fundadores Presentes:}
\vspace{-0.2cm}
\begin{enumerate}[noitemsep]
    \item Pardo Pozo Juan Diego (Jefe de Proyecto).
    \item Arnez Jaimes Mauricio.
    \item Quispe Flores Camila Belen.
    \item Ariñez Bautista Jim Gabriel.
    \item Medina Rodríguez Mateo Jhoustin.
    \item Molina Maldonado Tomas Manuel.
    \item Zeballos Crespo Jonathan.
    \item Gutierrez Guzmán Angheline.
\end{enumerate}

\newpage
\subsection{Desarrollo del Acta}

\noindent \textbf{Control de Asistencia del Team}\par
Se procedió a la toma de asistencia a la hora acordada, corroborando la presencia de los 8 socios fundadores de la grupo-empresa Byte Busters S.R.L. Al contar con el quórum necesario, se dio inicio formal a la sesión.

\vspace*{0.5cm}
\noindent \textbf{Lanzamiento de QA y Asignación de Roles por Células}\par
Se oficializó el \textit{Kick-off} de la fase de Aseguramiento de Calidad (QA) para el incremento actual. Para garantizar una revisión objetiva, se realizó una asignación cruzada de responsables de QA distribuidos por células de desarrollo, asegurando que ningún desarrollador pruebe su propio código y promoviendo la detección temprana de anomalías.

\vspace*{0.5cm}
\noindent \textbf{Actualización del \textit{Definition of Done} (DoD)}\par
El equipo acordó refinar el \textit{Definition of Done}. A partir de este momento, para que una Historia de Usuario se considere terminada, será requisito obligatorio adjuntar la evidencia de ejecución de los \textit{Test Cases} y, en caso de fallos, la correcta documentación de los \textit{Bug Reports} en la plataforma de gestión del proyecto.

\vspace*{0.5cm}
\noindent \textbf{Revisión Técnica de Impedimentos (Sync Extendido)}\par
Se llevó a cabo un espacio de sincronización extendido para abordar obstáculos técnicos. Se debatieron soluciones específicas para cuellos de botella identificados en el \textit{Backend} y la integración de la Base de Datos, mientras que en el \textit{Frontend} se revisaron dependencias visuales pendientes. Todos los impedimentos fueron abordados y se delegaron acciones concretas para su resolución.

\vspace*{0.5cm}
\noindent \textbf{Gestión del Sprint Backlog: Bloqueos y \textit{Carry Over}}\par
El \textit{Scrum Master} lideró la revisión del \textit{Sprint Backlog}. No se identificaron tareas bloqueadas por dependencias externas y todas las historias de usuario fueron finalizadas en el ciclo actual.

\vspace*{0.5cm}
\noindent \textbf{Alineación sobre el Repositorio de Documentación}\par
Se realizó un refinamiento técnico sobre las políticas del repositorio de documentación. Se establecieron convenciones claras para la redacción de READMEs, actualización de la wiki del proyecto y versionado de documentos técnicos, asegurando que la arquitectura y el código queden correctamente respaldados.

\vspace*{0.5cm}
\noindent \textbf{Aprobación Oficial}\par
Al no haber más puntos a tratar y con todas las dudas metodológicas y técnicas resueltas, se da por aprobada la sesión, validando los compromisos adquiridos en las áreas de desarrollo y QA.

\newpage
\noindent \textbf{Firmas y Aprobación de los Presentes}
\begin{center}
    \noindent
    \setlength{\tabcolsep}{0pt}
    \begin{tabular}{ccc}
        \espacioFirma{assets/img/firmas/mauricio.jpg}{Mauricio Jaimes Arnez}{13751221} & 
        \espacioFirma{assets/img/firmas/camila.jpg}{Camila Belen Quispe Flores}{13097244} &
        \espacioFirma{assets/img/firmas/jim.jpg}{Jim Gabriel Ariñez Bautista}{9517642} \\[1.0cm]
        \espacioFirma{assets/img/firmas/mateo.jpg}{Mateo Medina Rodríguez}{9453809} &
        \espacioFirma{assets/img/firmas/tomas.jpg}{Tomas Molina Maldonado}{8051701} &
        \espacioFirma{assets/img/firmas/jonathan.jpg}{Jonathan Zeballos Crespo}{13067494} \\[1.0cm]
    \end{tabular}
    \vspace{0.2cm} 
    \begin{tabular}{cc}
        \espacioFirma{assets/img/firmas/diego.jpg}{Juan Diego Pardo Pozo}{12747783} & 
        \espacioFirma{assets/img/firmas/angheline.jpg}{Angheline Gutierrez Guzmán}{13751815}
    \end{tabular}
\end{center}

\vspace{0.5cm}
\noindent \textbf{Fecha de última revisión:} Lunes, 20 de Abril de 2026